\begin{lemma}[Constructing and ordering the three exact Rutgers exits]
Let $\mathcal E$ be a fixed extended Tashkinov tree and $\mathcal P$ a
Rutgers closed presentation relative to $\mathcal E$, with coloring
$\sigma$ and displayed tree $U$. Suppose that $\sigma$ is proper and
$V(U)$ is elementary and closed under $\sigma$. Let $\alpha\ne\beta$
be colors such that $\alpha$ is missing on $U$ and $\alpha,\beta$ are
not $U$-interchangeable under $\sigma$. There is a Rutgers counterexample
on $\mathcal P$ with its two stored colors exactly $\alpha$ and $\beta$.
In particular its three exit packages retain individual boundary-edge
identities and have strictly increasing inside endpoints in the order of
$U$.
\end{lemma}
\begin{proof}
First let $c$ be any color missing on $U$. By \noderef{TreeMissingColors}
choose $v\in V(U)$ with $c$ missing at $v$. If an exact edge $f$ belonged
to the $c$-colored boundary, \noderef{BoundaryColorEdges} would give both
boundary membership and $\sigma(f)=c$. The assumed closedness, expanded
using \noderef{ColoredVertexSetClosed}, forbids this value when $c$ is
missing at $v$. Thus the $c$-colored boundary is empty. In particular
\noderef{TreeColorSetClosed} holds on the singleton $\{\alpha\}$.

If $\beta$ were also missing on $U$, the same argument would make its
colored boundary empty. Its cardinality would be zero, hence at most
one. Apply \noderef{ElementarySingleBoundaryInterchangeable} to
properness, the assumed elementarity, the empty $\alpha$-boundary, and
this cardinality bound. It would give $U$-interchangeability, including
isolated two-color path components, contrary to the hypothesis.
Therefore $\beta$ is not missing on $U$.

Use \noderef{ElementaryClosedNoninterchangeableExitPair} with the
proper coloring, elementary vertex set, distinct colors, missingness of
$\alpha$, singleton closedness of $\alpha$, and failed interchangeability.
It gives two walks $Q_1,Q_2$ and exact edges $f_1,f_2$. Each walk is an
exit with outside end $a_i$ and inside end $b_i$, and $f_i$ is a
$\beta$-colored boundary edge, incident with $b_i$, whose edge-node lies
on $Q_i$. Their supports are disjoint, and their outside ends and inside
ends are respectively distinct. Use
\noderef{ElementaryClosedExitPairExtendsTriple} with the same properness,
elementarity, distinct colors, missingness of $\alpha$, singleton
closedness, the proved nonmissingness of $\beta$, and these two exact
exit walks and disjointness and endpoint data. It produces a third
exit $Q_3$ and an exact edge $f_3$ with the same individual boundary,
incidence, and walk-membership properties. Its support is disjoint from
both preceding supports, and each of its ends is distinct from both
corresponding preceding ends. Thus all three outside ends, and all three
inside ends, are pairwise distinct.

Keep $(a_i,b_i,Q_i,f_i)$ as a single package for each $i$.
The inside-membership clause of \noderef{TreeExitPath} puts $b_i$ in
$V(U)$. By \noderef{TreeSequenceVertexSet}, select a displayed vertex
index $r_i$ with $U.\mathrm{vertex}(r_i)=b_i$. The injective vertex map
of \noderef{TreeSequence} makes each selected index unique. The pairwise
distinct inside ends give pairwise distinct indices. These are indices
in a linearly ordered finite interval, so exactly one of the following
six increasing orders holds:
\[
 (r_1,r_2,r_3),\ (r_1,r_3,r_2),\ (r_2,r_1,r_3),\
 (r_2,r_3,r_1),\ (r_3,r_1,r_2),\ (r_3,r_2,r_1).
\]
For completeness, compare $r_1$ with $r_2$. If $r_1<r_2$, comparison
of $r_2$ with $r_3$ gives the first order when $r_2<r_3$; otherwise
comparison of $r_1$ with $r_3$ gives the second or fifth order.
If $r_2<r_1$, comparison of $r_1$ with $r_3$ gives the third order
when $r_1<r_3$; otherwise comparison of $r_2$ with $r_3$ gives the
fourth or sixth order. Distinctness excludes equality in every comparison.

Relabel whole packages by the resulting order, never just one entry in
a package. Every individual exit, exact boundary membership, incidence,
and edge-node membership statement is unchanged. Pairwise distinctness
is preserved by a permutation. Pairwise support disjointness is preserved
too: the original assertion in the reverse order follows because
$x\in\operatorname{supp}(Q_j)$ and $x\in\operatorname{supp}(Q_i)$
would contradict the forward-order assertion. Finally the selected
display indices give $b_1\prec_U b_2\prec_U b_3$ exactly by
\noderef{TreeVertexOrder}. Fill all fields of
\noderef{RutgersCounterexample} with the presentation, the original
colors, their distinctness and missing/nonmissing facts, and the relabeled
packages and verified properties. The stored colors are $\alpha,\beta$
by construction, as claimed.
\end{proof}
